\documentclass[12pt,a4paper]{article}

% Packages
\usepackage[utf8]{inputenc}
\usepackage{amsmath,amssymb,amsthm}
\usepackage{mathrsfs}
\usepackage{geometry}
\usepackage[colorlinks=true, linkcolor=blue, citecolor=blue]{hyperref}
\usepackage{enumitem}
\usepackage{tikz-cd}
\usepackage{stmaryrd}
\usepackage{natbib}
\usepackage{fancyhdr}
\usepackage{lastpage}

\geometry{margin=1in}

% Theorem environments
\newtheorem{definition}{Definition}[section]
\newtheorem{theorem}[definition]{Theorem}
\newtheorem{proposition}[definition]{Proposition}
\newtheorem{lemma}[definition]{Lemma}
\newtheorem{corollary}[definition]{Corollary}
\newtheorem{remark}[definition]{Remark}
\newtheorem{example}[definition]{Example}

% Custom commands
\newcommand{\domain}{\mathcal{D}}
\newcommand{\vocab}{\mathcal{V}}
\newcommand{\lang}{\mathcal{L}}
\newcommand{\interp}{\mathcal{I}}
\newcommand{\concept}{\mathcal{C}}
\newcommand{\onto}{\mathcal{O}}
\newcommand{\intrel}{\mathfrak{R}}
\newcommand{\powerset}{\mathcal{P}}
\newcommand{\pfin}{\mathcal{P}_{\mathrm{fin}}}
\newcommand{\sem}[1]{\llbracket #1 \rrbracket}
\newcommand{\Phen}{\mathrm{Phen}}
\newcommand{\supp}{\mathrm{supp}}
% ===========================
% INFORMAÇÕES DO DOCUMENTO
% ===========================
\newcommand{\nomeautorum}{Igor Strozzi}
\newcommand{\nomeautordois}{Francesco Noseda}
\newcommand{\emailum}{istrozzi@matematica.ufrj.br}
\newcommand{\emaildois}{noseda@im.ufrj.br}
% Full title for title page
\newcommand{\ctitle}{Expanding Formal Ontologies \vspace{0.5cm}
	\\\Large{An Algebraic Framework}}
% Short title for header
\newcommand{\ctitleshort}{Expanding Formal Ontologies: An Algebraic Framework}
\newcommand{\disciplina}[1]{#1}	
\newcommand{\professor}{Francesco Noseda}
\newcommand{\universidade}{UFRJ}
\newcommand{\instituto}{Instituto de Matemática}
\newcommand{\departamento}{Departamento de Matemática Aplicada}

% ===========================
% CABEÇALHO E RODAPÉ
% ===========================
\pagestyle{fancy}
\fancyhf{}
%\lhead{\nomealuno}
\lhead{\ctitleshort}
%\chead{\ctitleshort}  % <-- Use short title here
%\rhead{\currentttl}
%\lfoot{\universidade}

%\cfoot{An algebraic framework}
\rfoot{Page \thepage\ of \pageref*{LastPage}}
\renewcommand{\headrulewidth}{0.4pt}
\renewcommand{\footrulewidth}{0.4pt}

% ===========================
% PÁGINA DE TÍTULO
% ===========================
\newcommand{\maketitlecustom}{
	\thispagestyle{empty}
	\begin{center}
		\hrule height 2pt width \textwidth
		\vspace{0.4cm}
		{\LARGE \textbf{\ctitle}} \\[0.4cm]
		\hrule height 2pt width \textwidth
		\vspace{0.8cm}
		{\Large \nomeautorum\footnote{\texttt{\emailum}} \quad \nomeautordois\footnote{\texttt{\emaildois}}} \\
		\vspace{0.8cm}
		\large{\textit{Applied Mathematics Department}\\
			Federal University of Rio de Janeiro}
	\end{center}
	\vspace{0.5cm}
	\hrule height 1pt
	\vspace{0.5cm}
}

%\title{Mathematical Foundations of Formal Ontologies\\
%\Large{Algebraic Structures for Ontological Expansion}}
%
%\author{Igor Strozzi\\
%Applied Mathematics Department\\
%Federal University of Rio de Janeiro}

\date{\today}

\begin{document}

\maketitlecustom


\begin{abstract}
\large{We explore the development of a mathematical foundation for formal ontologies within a purely set-theoretical framework, progressing from extensional relational structures through conceptualizations to ontology specifications. The central contribution lies in the articulation of \emph{ontological algebras} -- \textit{algebraic structures acting upon ordered ontologies that enable systematic ontological expansions while preserving structural coherence}. By grounding our framework in the minimal apparatus of partial orders and Heyting algebras, we intend to reveal the essential algebraic substrate underlying ontological reasoning. The resulting framework proposes to illuminate the dialectical relationship between static ontological representation and dynamic conceptual evolution.}
\end{abstract}

\newpage
\tableofcontents
\newpage
%=============================================================================
\section{Prolegomena: The Architecture of Formal Representation}
%=============================================================================

In the liminal space between abstract conceptualization and formal representation lies a profound epistemological challenge: how might one capture the invariant structures of a ``domain" while acknowledging the irreducible gap between symbolic articulation and conceptual content? The edifice of formal ontology emerges from this tension, offering systematic approaches to the formalization of domain knowledge while maintaining critical awareness of representation's inherent limitations.

The trajectory we trace here -- from extensional structures through conceptualizations to ontological algebras -- mirrors a fundamental dialectic in the theory of systems itself. Following \citet{gruber1993, guarino2009}, who characterized an ontology as ``an explicit specification of a conceptualization,'' we begin with the mathematical bedrock of sets and relations, enrich this substrate with conceptual structure, and ultimately arrive at algebraic mechanisms that enable ontologies to evolve while preserving their essential organization.

Our approach is distinguished by a commitment to minimalism: we work with the simplest ordered structures sufficient for our purposes (avoiding, for instance, any kind of explicit mereological or topological apparatus). This economy of means targets to reveal more clearly the essential algebraic character of ontological expansion.

The progression unfolds through four levels of increasing abstraction:
\begin{enumerate}[label=(\roman*)]
    \item \textbf{Extensional relational structures}: The mathematical substrate of sets and relations, providing the phenomenal cross-section of a domain
    \item \textbf{Intensional relational structures (Conceptualizations)}: Abstract, language-independent representations that transcend particular extensional instantiations by enriching them through an ``intensional turn"
    \item \textbf{Ontological commitments}: The semantic anchoring whereby formal languages latch onto conceptual structures
    \item \textbf{Ontological algebras}: The algebraic apparatus enabling compositional reasoning and systematic ontological expansion
\end{enumerate}

addressing related, relevant concepts as it feels needed.

%=============================================================================
\section{Introduction: of Ontology and ontologies}
%=============================================================================

A lengthier discussion on the difference between ``Ontology'' and ``(an) ontology" is done at \citet{guarino2009}. Not straying from what is discussed there, we try our own take here.

As a philosophical discipline, Ontology, oftentimes considered a subbranch of Metaphysics, aims to inquire on the most foundational and abstract aspects of ``what exists". It aims to study ``existent stuff", or ``Existence", or ``(the) Being", under the plain hypothesis of them \textit{being}, i.e, as \textit{existent entities}. Hence questions of extreme generality such as ``What does \textit{being} \textit{requires}? What does \textit{being} \textit{entails?}", which one could rephrase as ``which is the `smallest set of properties' a `thing' must have to be considered as `existent'?"; a question which might sound somewhat silly, but do ask yourself: does an abstract mathematical entity $M$ ``exist"? Why? 

Any answer should demand an ``ontological commitment'' from the entity who answers. Even further, it should in general require an ``ontological theory": some ``theory of Being" over which one could (reasonably?) rely to judge whether one such entity as $M$ exists or not.

Now, in the author's opinion, it is somewhat reassuring the fact that Ontology, as a broad realm of philosophical investigation, has found its way into more practical, concrete fields of study. If millenea of philosophical ontology inquiry suggests a paradigm that works within less general scopes, that can be actually a good paradigm to have. 

It is a strong idea, that of considering ``philosophical ontology with less general scopes". It actually allows us to leverage the whole ``conceptually-mature philosophical paradigm" we get from philosophical ontology to constrained scopes, or domains, what happens to be occasionally useful, to the point of earning a name. Let us call those ``domain-specific ontologies''. Pick your favorite physical theory -- which is a theory about reality. Think about what ``entities" such a theory imposes to exist, and about \textit{how} does it impose thoses entities to exist. Any physical theory imposes an ontological commitment. And the ontology that sublies the commitment is ``domain-specific'' inasmuch as one don't consider such theory to ``exhaust the whole of Being".

In a more practical setting, the problem is often stated the other way around. An applied ontologist has some \textit{modeling goal}, and needs to come up with a domain-specific ontology that's useful for their modeling goal. So there's no ``beforehand theory" about what exists, as in the physical theory exercise above. The process of ``engineering an ontology'' involves, precisely, jumping directly to a somewhat undefined ``domain of reality", cutting it off with what we came here to call a ``phenomenal cross-section of a domain", and bestowing that object with an intensional substrate. 

What is of most relevance to us here, in the discussions above, will be addressed more carefully and formally in the next sections.

%=============================================================================
\section{Preliminary Definitions}
%=============================================================================

Although we, in general, accept the definitions found in \citet{guarino2009} we adopt a slightly different order of exposition, by choosing to discuss and present (even though informally, on occasions) notions such as that of a \textit{system} and of a \textit{domain} in advance. Hopefully, the motivations for the following definitions will be clear in due time.

Let us first discuss the concept of a system, similarly as presented in \citet{bertalanffy1968}.

A system can be defined as a ``complex of interacting elements", where by interactions we can understand that for some relation $R$ (of arity $n$), defined over the ``set of elements" $P$, for some $(p_i)_{i \in n} \in P^n$, $(p_i)_{i \in n} \in R$ is relevant for the system's \textit{global} behavior: that is what captures ``complex" in the ``interacting of things". With that considered, a familiar-to-be mathematical object is suggested: that of an extensional relational structure.

See, if a single relation with a single element suffices, as long as it impacts the system's behavior as a whole, to be considered in the definition of a system (as a ``complex of interacting elements"), it is only natural to consider that a system would require in its definition possibly many relations, with, why not, many different arities. It would appeal to reason, then, to consider:

\begin{definition}[System (Preliminary)]
\label{def:system-prelim}
	Let $X$ be a non-empty set. Let $R^{\langle i \rangle}_j$ be an $i$-ary relation over $X$. A \textit{system} is a tuple $\mathcal{S} = (X, R)$, where $R = (\{R^{\langle i \rangle}_j\}_{j\in J_i})_{i\in \mathbb{Z}_+}$ for some index sets $J_i$.
\end{definition}

This preliminary definition, while serviceable, obscures an important structural feature: the \emph{grading by arity}. The superscript notation $R^{\langle i \rangle}$ hints at an organizational principle---relations are stratified by their arity---but the notation treats this stratification as mere bookkeeping rather than constitutive structure. We shall shortly provide a reformulation that elevates this grading to definitional status.

%=============================================================================
\subsection{The Graded Relational Space}
%=============================================================================

To make the arity-grading explicit and to prepare the ground for the notion of \emph{phenomenal cross-section}, we introduce the graded relational space over a domain.

\begin{definition}[Graded Relational Space]
\label{def:graded-rel-space}
Let $\domain$ be a non-empty set (the \emph{domain of discourse}). The \emph{graded relational space} over $\domain$ is the family $E = (E_n)_{n \in \mathbb{Z}_+}$ where
$$
E_n = \{R \subseteq D_1 \times \cdots \times D_n \mid D_j \in \powerset(\domain)\}
$$
is the collection of all $n$-ary relations over (possibly heterogeneous) subsets of $\domain$.
\end{definition}

A word on the definition is in order. By allowing relations to be defined over subsets $D_j \subseteq \domain$ rather than over $\domain$ itself, we accommodate the informal notion of \emph{sorts} or \emph{types} within the domain. Elements of $\domain$ may be heterogeneous---persons, events, abstract entities---and a relation like ``reports-to" naturally restricts its arguments to certain kinds. We do not formalize a sort system here; the powerset construction $D_j \in \powerset(\domain)$ suffices to capture this flexibility without imposing additional apparatus.

The fiber $E_n$ collects all possible $n$-ary relations---not merely those that happen to obtain in some particular state of affairs, but all those that \emph{could} meaningfully be predicated of $n$-tuples drawn from the domain. This totality is, of course, vast; the systems and worlds we define will select finite portions of it.

%=============================================================================
\subsection{Multi-Sections}
%=============================================================================

The graded relational space provides the ambient structure; we now require a notion of \emph{selection} from this space that respects the arity grading.

\begin{definition}[Multi-Section]
\label{def:multi-section}
A \emph{multi-section} over $\domain$ is a function $s : \mathbb{Z}_+ \to \pfin(E)$ satisfying:
\begin{enumerate}[label=(\roman*)]
    \item $s(n) \subseteq E_n$ for all $n$ \hfill (arity-respecting)
    \item $\supp(s) := \{n \in \mathbb{Z}_+ \mid s(n) \neq \emptyset\}$ is finite \hfill (finite support)
\end{enumerate}
The \emph{extent} of $s$ is $|s| = \bigcup_{n \in \mathbb{Z}_+} s(n)$.
\end{definition}

The terminology ``multi-section" is chosen advisedly. In the theory of fiber bundles, a section selects exactly one element from each fiber; here, we permit \emph{multiple} selections at each arity (hence ``multi"), and we allow empty selection (no relations of a given arity need be present). The finite support condition ensures that our structures remain finitary---a system involves only finitely many arities, even if the domain itself is infinite.

The extent $|s|$ ``flattens" the graded structure into a mere set of relations, discarding the arity information. This flattening is useful when passing to the classical notion of extensional relational structure, but the graded structure $s$ itself carries more information: it tells us not only \emph{which} relations obtain, but \emph{how many relations of each arity} the structure involves.

%=============================================================================
\subsection{Systems Reformulated}
%=============================================================================

With the apparatus of multi-sections in hand, we can now provide a cleaner formulation of the notion of system.

\begin{definition}[System]
\label{def:system}
A \emph{system} over $\domain$ is a multi-section $S : \mathbb{Z}_+ \to \pfin(E)$ over the graded relational space of $\domain$.

We write $\mathcal{S} = (\domain, S)$ when emphasizing the domain, and identify $\mathcal{S}$ with the pair $(\domain, |S|)$ when passing to the associated extensional relational structure.
\end{definition}

\begin{remark}[Equivalence with Preliminary Definition]
This reformulates Definition~\ref{def:system-prelim}: the tuple $(X, (\{R^{\langle i \rangle}_j\}_{j \in J_i})_{i \in \mathbb{Z}_+})$ corresponds to a multi-section where $S(i) = \{R^{\langle i \rangle}_j \mid j \in J_i\}$ for finite index sets $J_i$. The grading by arity---implicit in the superscript notation of the preliminary definition---becomes the defining structural feature in the reformulation.
\end{remark}

The virtue of this reformulation lies not in increased generality but in increased \emph{perspicuity}. By making the graded structure explicit, we prepare the ground for the crucial notion of phenomenal cross-section: a \emph{partial view} of a system, obtained by selecting some but not all of its relations.

%=============================================================================
\subsection{Worlds}
%=============================================================================

Before introducing phenomenal cross-sections, we pause to define the notion of a \emph{world}---a complete specification of which relations obtain at a given state. This notion bridges our framework to that of \citet{guarino2009}, where worlds (or ``possible worlds") play a central role in the intensional account of meaning.

\begin{definition}[World]
\label{def:world}
Let $\domain$ be a domain. A \emph{world} over $\domain$ is a multi-section $w : \mathbb{Z}_+ \to \pfin(E)$.
\end{definition}

One might object: this is identical to the definition of a system! Indeed, at the level of mathematical structure, a world and a system are the same kind of object---a multi-section over the graded relational space. The distinction is one of \emph{interpretation} and \emph{role within the theory}:

\begin{itemize}
    \item A \textbf{system} represents a ``piece of reality" we wish to model---a company, an ecosystem, a social network. It is the target of ontological investigation.
    \item A \textbf{world} represents a ``maximal observable state of affairs"---a complete specification of which relations hold at a given moment or under a given configuration.
\end{itemize}

The terminology follows \citet{guarino2009}, who characterize a world as ``a unique assignment of values to all the observable variables that characterize the system." In our framework, the ``observable variables" are precisely the relations of various arities, and a world specifies, for each arity $n$, which $n$-ary relations obtain.

When we consider a system evolving over time, or varying across possible configurations, we obtain a \emph{family} of worlds. This leads to the notion of domain space.

%=============================================================================
\subsection{Domain Spaces}
%=============================================================================

\begin{definition}[Domain Space]
\label{def:domain-space}
A \emph{domain space} is a pair $\langle \domain, W \rangle$ where $\domain$ is a domain and $W$ is a set of worlds over $\domain$.
\end{definition}

This definition corresponds precisely to that of \citet{guarino2009}, who write: ``The pair $\langle D, W \rangle$ is called a domain space for $S$, as it intuitively fixes the space of variability of the universe of discourse $D$ with respect to the possible states of $S$."

The set $W$ represents the \emph{admissible configurations}---those states of affairs deemed possible for the system under consideration. Not every mathematically conceivable multi-section need belong to $W$; the domain space encodes constraints on what is possible. A chemical system, for instance, might exclude worlds violating conservation laws; a social system might exclude worlds where the same person simultaneously occupies incompatible roles.

The domain space is the locus where \emph{extensional} and \emph{intensional} accounts of meaning meet. A particular world $w \in W$ provides an extensional snapshot; the totality $W$ provides the space over which intensional (conceptual) relations are defined.

%=============================================================================
\subsection{Phenomenal Cross-Sections}
%=============================================================================

We arrive at the central notion motivating our reformulation: the \emph{phenomenal cross-section}, which captures the idea of a \emph{partial view} of a system.

\begin{definition}[Phenomenal Cross-Section]
\label{def:phenom-cross-section}
Let $S$ be a system (i.e., a multi-section). A \emph{phenomenal cross-section} of $S$ is a multi-section $\sigma$ such that
$$
\sigma(n) \subseteq S(n) \quad \text{for all } n \in \mathbb{Z}_+
$$
We write $\sigma \sqsubseteq S$ and denote the collection of all phenomenal cross-sections of $S$ by $\Phen(S)$.
\end{definition}

The terminology ``phenomenal cross-section" is chosen to evoke the epistemological situation of the ontologist. A system---a company, say, with all its employees, hierarchies, collaborations, and dependencies---exists as a complex totality. But no ontology captures this totality in full; every ontology commits to certain relations (``reports-to," ``cooperates-with") while remaining silent on others (perhaps ``is-taller-than," or ``shares-birthday-with"). The phenomenal cross-section $\sigma$ represents what is \emph{visible} or \emph{accessible} to a given ontological commitment: the relations that the ontology acknowledges as meaningful.

The gap $S \setminus \sigma$---those relations present in the system but absent from the cross-section---constitutes the ontologically ``invisible" structure: relations that obtain but that the ontology does not represent. This gap is not a failure of the ontology but rather a constitutive feature of any finite specification before the richness of the real.

\begin{proposition}[Lattice Structure]
\label{prop:phen-lattice}
$(\Phen(S), \sqsubseteq)$ is a bounded distributive lattice with:
\begin{itemize}
    \item $\bot(n) = \emptyset$ \hfill (empty observation)
    \item $\top = S$ \hfill (complete observation)
    \item $(\sigma \wedge \tau)(n) = \sigma(n) \cap \tau(n)$ \hfill (meet)
    \item $(\sigma \vee \tau)(n) = \sigma(n) \cup \tau(n)$ \hfill (join)
\end{itemize}
Indeed, $\Phen(S) \cong \prod_{n} \powerset(S(n))$, a product of Boolean algebras.
\end{proposition}

\begin{proof}
The isomorphism $\Phen(S) \cong \prod_n \powerset(S(n))$ is immediate: a phenomenal cross-section $\sigma$ is determined by choosing, for each $n$, a subset of $S(n)$. The lattice operations are inherited pointwise from the Boolean algebra structure on each $\powerset(S(n))$.
\end{proof}

The lattice structure on $\Phen(S)$ admits a natural interpretation:
\begin{itemize}
    \item \textbf{Meet} $\sigma \wedge \tau$: the relations acknowledged by \emph{both} cross-sections---the common observations.
    \item \textbf{Join} $\sigma \vee \tau$: the relations acknowledged by \emph{either} cross-section---the pooled observations.
    \item \textbf{Order} $\sigma \sqsubseteq \tau$: $\tau$ is a \emph{refinement} of $\sigma$, acknowledging everything $\sigma$ does and possibly more.
\end{itemize}

Different phenomenal cross-sections thus represent different \emph{ontological commitments} or \emph{observational capacities}. The lattice structure captures how these commitments can be compared, combined, and refined.

%=============================================================================
\section{Extensional Foundations}
%=============================================================================

We now connect our graded framework to the classical notion of extensional relational structure, which provides what we have termed the phenomenal cross-section of a domain.

\begin{definition}[Extensional Relational Structure]
An \emph{extensional relational structure} \citep[cf.][]{genesereth1987} is a tuple $\mathcal{S} = (\domain, \mathbf{R})$ where:
\begin{itemize}
    \item $\domain$ is a non-empty set called the \emph{universe of discourse}
    \item $\mathbf{R} = \{R_i\}_{i \in I}$ is a family of relations on $\domain$, where each $R_i \subseteq \domain^{n_i}$ for some arity $n_i \in \mathbb{N}$
\end{itemize}
\end{definition}

\begin{remark}[Relation to Multi-Sections]
Every multi-section $s$ over $\domain$ determines an extensional relational structure $(\domain, |s|)$ by taking the extent. Conversely, every extensional relational structure $(\domain, \mathbf{R})$ with finite $\mathbf{R}$ determines a multi-section $s$ by setting $s(n) = \{R \in \mathbf{R} \mid R \text{ is } n\text{-ary}\}$. The multi-section carries strictly more information: it records the arity of each relation as part of the structure, rather than as a derived property.
\end{remark}

This deceptively simple mathematical structure conceals profound ontological choices. The extensional relational structure captures entities and their actual interrelations at a particular state---a snapshot devoid of modal or intensional content. Yet this very extensionality creates a constitutive tension: the structure cannot inherently distinguish between relations that are intrinsic to the system's organization and relations that merely happen to obtain.

\begin{remark}[The Constitutive-Summative Distinction]
Following \citeauthor{bertalanffy1968}'s foundational insight \citep{bertalanffy1968}, we recognize that an extensional structure $(\domain, \mathbf{R})$ captures relations as sets of tuples but cannot discriminate between:
\begin{enumerate}
    \item Relations that are \emph{constitutive}---intrinsic to the system's identity and organization
    \item Relations that are \emph{summative}---contingent aggregations without organizational significance
\end{enumerate}
This limitation is not a deficiency to be overcome but rather a constitutive feature of extensional representation itself, one that motivates the ascent to conceptualization.
\end{remark}

\begin{example}[Isomeric Structures]
Consider chemical isomers sharing identical molecular formulas yet exhibiting distinct structural arrangements. Within the extensional framework:
\begin{align*}
    \domain &= \{\text{atom}_1, \text{atom}_2, \ldots, \text{atom}_n\}\\
    \mathbf{R}_{\alpha} &= \{\text{bonds-to}_\alpha, \text{orbital-overlap}_\alpha, \ldots\}\\
    \mathbf{R}_{\beta} &= \{\text{bonds-to}_\beta, \text{orbital-overlap}_\beta, \ldots\}
\end{align*}
The emergent chemical properties---divergent reactivities, distinct spectroscopic signatures---exemplify constitutive characteristics arising from specific relational configurations. Yet the extensional structure alone provides no principle for distinguishing these constitutive relations from mere accidental correlations.
\end{example}

%=============================================================================
\section{Conceptualization}
%=============================================================================

The conceptualization represents our ascent from the extensional plane to the intensional---an abstract, language-independent representation of a domain that serves as the semantic foundation for ontological specification \citep{gruber1995, guarino1998}.

%-----------------------------------------------------------------------------
\subsection{From Worlds to Conceptual Relations}
%-----------------------------------------------------------------------------

The key insight bridging our framework to the intensional account is this: a \emph{conceptual relation} is not a single extensional relation but a \emph{function} assigning an extensional relation to each world. This is the ``intensional turn" that distinguishes ontological engineering from mere database design.

Recall that a domain space $\langle \domain, W \rangle$ fixes the domain and the set of admissible worlds. Given such a domain space, we can define:

\begin{definition}[Conceptual Relation]
\label{def:conceptual-relation}
Let $\langle \domain, W \rangle$ be a domain space. A \emph{conceptual relation} (or \emph{intensional relation}) of arity $n$ on $\langle \domain, W \rangle$ is a total function
$$
\rho^n : W \to E_n
$$
from the set of worlds into the fiber of $n$-ary relations.
\end{definition}

This definition captures the essential idea: the \emph{meaning} of a conceptual relation like ``cooperates-with" is not exhausted by its extension in any particular world (who cooperates with whom today) but is rather given by the \emph{rule} that determines its extension in \emph{every} admissible world. Two agents share the same concept of ``cooperation" if, presented with any world state, they would agree on which pairs of individuals are cooperating.

The extensional content at a particular world $w$ is recovered as $\rho^n(w) \in E_n$---the $n$-ary relation that ``obtains" for that conceptual relation in world $w$.

%-----------------------------------------------------------------------------
\subsection{Conceptualization Defined}
%-----------------------------------------------------------------------------

\begin{definition}[Conceptualization]
\label{def:conceptualization}
A \emph{conceptualization} is a triple $\concept = (\domain, \Sigma, \intrel)$ where:
\begin{itemize}
    \item $\domain$ is the universe of discourse
    \item $\Sigma$ is a set of \emph{conceptual signatures} specifying the types and arities of meaningful relations
    \item $\intrel = \{\rho_\sigma\}_{\sigma \in \Sigma}$ is a family of \emph{conceptual relations}, where each $\rho_\sigma$ is a conceptual relation of arity $\text{arity}(\sigma)$ on some domain space $\langle \domain, W \rangle$
\end{itemize}
\end{definition}

\begin{remark}[Implicit Domain Space]
The definition of conceptualization presupposes a domain space $\langle \domain, W \rangle$ over which the conceptual relations are defined. This domain space is often left implicit, with $W$ understood as ``all possible worlds consistent with our conceptual understanding." We make it explicit here to clarify the connection with our framework of systems and worlds.
\end{remark}

The conceptualization transcends mere extensional enumeration by incorporating \emph{intensional} content through its signature structure. The set $\Sigma$ specifies not merely what relations exist, but what \emph{kinds} of relations are meaningful within the domain---a crucial distinction that separates ontological engineering from mere database design.

%-----------------------------------------------------------------------------
\subsection{Intended Models}
%-----------------------------------------------------------------------------

\begin{definition}[Intended Models]
Given a conceptualization $\concept = (\domain, \Sigma, \intrel)$ over a domain space $\langle \domain, W \rangle$, the set of \emph{intended models} $I(\concept)$ consists of all extensional relational structures $\mathcal{S} = (\domain, \mathbf{R})$ satisfying \citep[following][]{guarino2009}:
\begin{enumerate}
    \item The domain of $\mathcal{S}$ coincides with that of $\concept$
    \item There exists a world $w \in W$ such that for each $\sigma \in \Sigma$, there exists $R \in \mathbf{R}$ with $R = \rho_\sigma(w)$
\end{enumerate}
\end{definition}

The intended models constitute the extensional ``shadow'' cast by a conceptualization---the space of possible states of affairs consistent with our conceptual understanding. They represent the bridge between intensional content and extensional realization.

\begin{remark}[Intended Models as Phenomenal Cross-Sections]
In our framework, an intended model corresponds to a phenomenal cross-section of a particular world. Given a world $w \in W$ and a conceptualization $\concept$, the extensional structure $(\domain, \{\rho_\sigma(w)\}_{\sigma \in \Sigma})$ is precisely the cross-section of $w$ that ``sees" only those relations named by the signature $\Sigma$. The conceptualization thus acts as a \emph{filter}, selecting which relations are ontologically relevant.
\end{remark}

%=============================================================================
\section{Ontological Commitment and Specification}
%=============================================================================

\subsection{The Nature of Commitment}

Ontological commitment performs the crucial mediating work of anchoring vocabulary to conceptualization, specifying how linguistic terms latch onto conceptual structure. The notion traces to \citeauthor{quine1948}'s dictum that ``to be is to be the value of a bound variable'' \citep{quine1948}, though here we extend it to intensional contexts following \citet{guarino1998}.

\begin{definition}[Ontological Commitment]
Let $\lang$ be a first-order logical language with vocabulary $\vocab$ and let $\concept = (\domain, \Sigma, \intrel)$ be a conceptualization. An \emph{ontological commitment} for $\lang$ is a tuple $K = (\concept, \interp)$ where the \emph{interpretation function} $\interp : \vocab \to \domain \cup \intrel$ is a total function mapping:
\begin{itemize}
    \item Each constant symbol in $\vocab$ to an element of $\domain$
    \item Each predicate symbol in $\vocab$ to a conceptual relation in $\intrel$
\end{itemize}
\end{definition}

The term ``commitment'' carries profound philosophical weight. To commit to a conceptualization $\concept$ is to specify how one's vocabulary latches onto the conceptual structure of $\concept$---not merely onto some particular extensional snapshot, but onto the relations that constitute the domain's organization. This addresses a deep problem in ontological engineering: without such anchoring, the same linguistic theory could be satisfied by radically different conceptualizations \citep{guarino2009}.

\begin{remark}[Commitment and Phenomenal Cross-Section]
The ontological commitment $K = (\concept, \interp)$ determines which conceptual relations are ``visible" to the language $\lang$: precisely those in the image of $\interp$. If $\concept$ contains conceptual relations not named by any predicate in $\vocab$, these remain ontologically silent---present in the conceptualization but inexpressible in the language. The commitment thus induces a phenomenal cross-section of the conceptualization itself.
\end{remark}

\begin{definition}[Compatibility and Intended Models]
An extensional first-order structure $M = (S, I)$ for $\lang$ is \emph{compatible with} ontological commitment $K = (\concept, \interp)$ if:
\begin{enumerate}
    \item For all constant symbols $c \in \vocab$: $I(c) = \interp(c)$
    \item There exists a world $w$ such that for each predicate symbol $p \in \vocab$: $I(p) = \interp(p)(w)$
\end{enumerate}
The set $I_K(\lang)$ of all compatible models constitutes the \emph{intended models of $\lang$ according to $K$}.
\end{definition}

\subsection{Ontology as Approximation}

\begin{definition}[Ontology]
An \emph{ontology} for a conceptualization $\concept$ with vocabulary $\vocab$ and ontological commitment $K$ is a logical theory $\onto$ consisting of a set of formulas of $\lang$, designed so that the set of its models approximates the set of intended models $I_K(\lang)$ \citep{gruber1993, guarino1998}.
\end{definition}

The language of ``approximation'' here is not accidental but rather constitutive. An ontology does not \emph{fully capture} a conceptualization but rather constrains the space of models to approach the intended models. The gap between $\text{Mod}(\onto)$ and $I_K(\lang)$ measures our axiomatic incompleteness---a gap that is not mere deficiency but rather an inevitable feature of formal representation before the richness of conceptual content.

%=============================================================================
\section{Ordered Ontologies}
%=============================================================================

We now introduce the minimal ordered structure sufficient for algebraic action. In the interstices between arbitrary set-theoretic structure and rich topological apparatus, we discover that a simple partial order provides the essential substrate for ontological algebra \citep[for foundational material on order theory, see][]{davey2002}.

\begin{definition}[Ordered Ontology]
An \emph{ordered ontology} is a pair $\onto = (O, \leq)$ where:
\begin{itemize}
    \item $O$ is a non-empty set of \emph{ontological entities} (concepts, individuals, or relations)
    \item $\leq$ is a partial order on $O$, satisfying:
    \begin{enumerate}[label=(\alph*)]
        \item \textbf{Reflexivity}: $\forall x \in O.\; x \leq x$
        \item \textbf{Transitivity}: $x \leq y \land y \leq z \implies x \leq z$
        \item \textbf{Antisymmetry}: $x \leq y \land y \leq x \implies x = y$
    \end{enumerate}
\end{itemize}
\end{definition}

The partial order $\leq$ admits multiple interpretations depending on the ontological domain: subsumption in taxonomic hierarchies, specificity in concept lattices, or logical strength in theories. This interpretive flexibility is a virtue---the algebraic machinery we develop operates uniformly across these diverse readings.

\begin{definition}[Derived Relations]
From the primitive order, we define:
\begin{align*}
    x < y &:\iff x \leq y \land x \neq y & \text{(strict subordination)}\\
    x \parallel y &:\iff \neg(x \leq y) \land \neg(y \leq x) & \text{(incomparability)}\\
    x \sim y &:\iff x \leq y \lor y \leq x & \text{(comparability)}
\end{align*}
\end{definition}

\begin{remark}[The Lattice Completion]
Given an ordered ontology $(O, \leq)$, one may always consider its Dedekind-MacNeille completion---the smallest complete lattice containing $(O, \leq)$ as a sub-poset \citep{davey2002}. This completion provides meets and joins for arbitrary subsets, though we shall not require this full apparatus for our development.
\end{remark}

%=============================================================================
\section{Compositional Algebras}
%=============================================================================

The Heyting algebra emerges as the natural algebraic structure for compositional ontological reasoning. Its operations encode the fundamental modes of conceptual combination, while its intrinsic constructive character reflects the epistemological constraints inherent in ontological work \citep{heyting1956}.

\subsection{The Heyting Structure}

\begin{definition}[Heyting Algebra]
A \emph{Heyting algebra} \citep{heyting1956, johnstone1982} is a bounded distributive lattice $H = (H, \sqcup, \sqcap, \Rightarrow, 0, 1)$ such that for all $a, b \in H$, the \emph{relative pseudo-complement} $a \Rightarrow b$ exists, characterized by the adjunction:
$$
c \sqcap a \leq b \iff c \leq a \Rightarrow b
$$
\end{definition}

This adjunction property---the defining characteristic of Heyting algebras---encodes a profound correspondence between conjunction and implication. It asserts that asking ``when does $c$ combined with $a$ yield something below $b$?'' is equivalent to asking ``when is $c$ below the conditional $a \Rightarrow b$?'' This duality between combination and entailment pervades compositional reasoning.

\begin{proposition}[Fundamental Properties]
In any Heyting algebra:
\begin{enumerate}[label=(\roman*)]
    \item The negation $\sim a := a \Rightarrow 0$ satisfies $a \sqcap \sim a = 0$
    \item Contraposition holds: $a \leq b \implies \sim b \leq \sim a$
    \item $a \leq \sim\sim a$ but generally $\sim\sim a \not\leq a$
    \item Modus ponens: $a \sqcap (a \Rightarrow b) \leq b$
\end{enumerate}
\end{proposition}

The failure of double negation elimination---property (iii)---reflects the constructive character of Heyting algebras. One cannot, in general, recover $a$ from the mere absence of its negation; positive construction is required. This constructive orientation aligns naturally with the epistemological situation in ontological engineering, where one cannot assert the existence of concepts without constructive justification \citep[cf.][]{sorensen2006}.

\subsection{Compositional Interpretation}

The operations of a Heyting algebra admit natural interpretation as modes of conceptual composition:

\begin{center}
\renewcommand{\arraystretch}{1.3}
\begin{tabular}{lll}
\textbf{Operation} & \textbf{Notation} & \textbf{Conceptual Reading}\\
\hline
Meet & $a \sqcap b$ & Conjunction; the composite concept\\
Join & $a \sqcup b$ & Disjunction; conceptual choice\\
Implication & $a \Rightarrow b$ & Entailment; conceptual dependency\\
Negation & $\sim a$ & Complement; conceptual exclusion\\
Top & $1$ & Universal concept; maximal generality\\
Bottom & $0$ & Impossible concept; contradiction
\end{tabular}
\end{center}

\begin{definition}[Algebraic Derivability]
In a Heyting algebra $A$, the \emph{derivability relation} is defined:
$$
a \vdash_A b :\iff a \sqcap b = a
$$
This is equivalent to $a \leq b$ in the underlying order.
\end{definition}

Derivability captures the sense in which one concept ``contains'' or ``entails'' another within the algebraic structure. The equivalence $a \vdash_A b \iff a \leq b$ reveals that algebraic derivability precisely mirrors the order-theoretic subsumption.

%=============================================================================
\section{The Ontological Algebra}
%=============================================================================

We arrive at the central construction: the ontological algebra, wherein an algebraic structure acts upon an ordered ontology, enabling both compositional reasoning and systematic expansion.

\begin{definition}[Ontological Algebra]
An \emph{ontological algebra} is a triple $(\onto, A, \sem{\cdot})$ where:
\begin{itemize}
    \item $\onto = (O, \leq)$ is an ordered ontology
    \item $A = (A, \sqcup, \sqcap, \Rightarrow, 0, 1)$ is a Heyting algebra
    \item $\sem{\cdot} : O \to A$ is the \emph{interpretation function}
\end{itemize}
satisfying the \textbf{monotonicity condition}:
$$
c_1 \leq c_2 \implies \sem{c_1} \leq \sem{c_2}
$$
\end{definition}

The interpretation function $\sem{\cdot}$ performs the crucial work of embedding the ordered ontology into the richer algebraic structure. Monotonicity ensures that the ontological order is preserved under this embedding---subordination in the ontology translates to subordination in the algebra.

\begin{definition}[Induced Operations]
Given ontological entities $c_1, c_2 \in O$, we define:
\begin{align*}
    c_1 \triangleright c_2 &:= \sem{c_1} \sqcap \sem{c_2} \in A & \text{(combination)}\\
    c_1 \multimap c_2 &:= \sem{c_1} \Rightarrow \sem{c_2} \in A & \text{(entailment)}
\end{align*}
\end{definition}

These induced operations lift the algebraic structure to the ontological level. The combination $c_1 \triangleright c_2$ represents the ``joint content'' of two concepts as computed in the algebra, while the entailment $c_1 \multimap c_2$ captures the degree to which $c_1$ implies or grounds $c_2$.

\subsection{Fundamental Theorems}

\begin{theorem}[Order Preservation]
If $c_1 \leq c_2$ in $\onto$, then $\sem{c_1} \vdash_A \sem{c_2}$.
\end{theorem}

\begin{proof}
By monotonicity, $c_1 \leq c_2$ implies $\sem{c_1} \leq \sem{c_2}$. In a Heyting algebra, $a \leq b$ is equivalent to $a \sqcap b = a$, hence $\sem{c_1} \sqcap \sem{c_2} = \sem{c_1}$, which is the definition of $\sem{c_1} \vdash_A \sem{c_2}$.
\end{proof}

\begin{theorem}[Modus Ponens for Concepts]
For any $c_1, c_2 \in O$:
$$
\sem{c_1} \sqcap (c_1 \multimap c_2) \leq \sem{c_2}
$$
\end{theorem}

\begin{proof}
This is the fundamental Heyting algebra property: $a \sqcap (a \Rightarrow b) \leq b$ for all $a, b$.
\end{proof}

\begin{theorem}[Combination Bounds]
For any $c_1, c_2 \in O$:
$$
c_1 \triangleright c_2 \leq \sem{c_1} \quad \text{and} \quad c_1 \triangleright c_2 \leq \sem{c_2}
$$
\end{theorem}

\begin{proof}
Immediate from the lattice property: $a \sqcap b \leq a$ and $a \sqcap b \leq b$.
\end{proof}

%=============================================================================
\section{Ontological Expansion}
%=============================================================================

We now address the central question animating this investigation: how might a formal ontology evolve while maintaining structural coherence? The ontological algebra provides the mechanism through \emph{expansion operators}---algebraic actions that systematically extend the ontological domain.

\subsection{The Expansion Framework}

\begin{definition}[Ontological Expansion]
Let $(\onto, A, \sem{\cdot})$ be an ontological algebra with $\onto = (O, \leq)$. An \emph{expansion} is a tuple $(O', \leq', \sem{\cdot}')$ where:
\begin{enumerate}
    \item $O' \supseteq O$ (domain extension)
    \item $\leq'$ is a partial order on $O'$ extending $\leq$ (order extension)
    \item $\sem{\cdot}' : O' \to A$ extends $\sem{\cdot}$, satisfying monotonicity on $(O', \leq')$
\end{enumerate}
The expansion is \emph{conservative} if $\leq' \cap (O \times O) = \leq$.
\end{definition}

Conservative expansion preserves all existing order relations while potentially adding new entities and new relations involving those entities. This conservativity ensures that established ontological commitments remain intact under expansion.

\subsection{Algebraically Generated Expansion}

The key insight enabling systematic expansion is that elements of the Heyting algebra $A$ can serve as \emph{expansion parameters}---specifications of how new concepts relate to existing ones.

\begin{definition}[Concept Generation]
Given $c \in O$ and $a \in A$ with $a \leq \sem{c}$, the \emph{generated concept} $c \otimes a$ is defined by:
$$
\sem{c \otimes a}' := a
$$
with order relations:
$$
c \otimes a \leq' c
$$
and for any $c' \in O$:
$$
c \otimes a \leq' c' \iff a \leq \sem{c'}
$$
\end{definition}

The condition $a \leq \sem{c}$ ensures that generated concepts are subordinate to their generators---we create refinements or specializations, not arbitrary additions. The element $a \in A$ specifies precisely ``how much'' of $c$ the new concept captures.

\begin{proposition}[Well-Definedness of Generation]
The generated concept $c \otimes a$ yields a valid expansion: the extended structure $(O', \leq', \sem{\cdot}')$ satisfies monotonicity.
\end{proposition}

\begin{proof}
We must verify that $x \leq' y$ implies $\sem{x}' \leq \sem{y}'$. For $x, y \in O$, this follows from the original monotonicity. For $x = c \otimes a$ and $y \in O$, we have $c \otimes a \leq' y$ iff $a \leq \sem{y}$, and $\sem{c \otimes a}' = a$, so $\sem{c \otimes a}' \leq \sem{y}' = \sem{y}$ as required.
\end{proof}

\subsection{The Expansion Monoid}

\begin{definition}[Expansion Operator]
An \emph{expansion operator} is a function $\mathcal{E} : O \times A \to O'$ satisfying:
\begin{enumerate}
    \item $\mathcal{E}(c, \sem{c}) = c$ (identity at interpretation)
    \item $\mathcal{E}(c, a)$ is defined only when $a \leq \sem{c}$ (subordination constraint)
    \item The result extends the ontological algebra structure
\end{enumerate}
\end{definition}

\begin{proposition}[Monoid Structure]
The set of expansion operators forms a monoid under composition:
\begin{itemize}
    \item \textbf{Composition}: $(\mathcal{E}_1 \circ \mathcal{E}_2)(c, a) := \mathcal{E}_1(\mathcal{E}_2(c, a'), a'')$ where $a = a' \sqcap a''$
    \item \textbf{Identity}: $\text{id}(c, a) := c$ when $a = \sem{c}$
\end{itemize}
\end{proposition}

This monoid structure captures the compositionality of ontological expansion: sequential refinements can be combined, and ``doing nothing'' (expanding by the interpretation itself) serves as identity.

\subsection{Semantic Characterization}

\begin{theorem}[Expansion as Filtering]
Let $(\onto, A, \sem{\cdot})$ be an ontological algebra and $(O', \leq', \sem{\cdot}')$ an expansion generated by $c \otimes a$. Then for any formula $\phi$ in the language of $\onto$:
$$
\onto' \models \phi[c \otimes a / x] \iff \onto \models \phi[c/x] \text{ and } a \text{ satisfies the constraints encoded by } \phi
$$
\end{theorem}

This theorem reveals expansion as a \emph{filtering} operation: the generated concept $c \otimes a$ inherits properties from $c$ but only those compatible with the algebraic constraint $a$. Expansion thus provides controlled specialization within the algebraic framework.

%=============================================================================
\section{The Dialectics of Formalization}
%=============================================================================

The trajectory from extensional structures through conceptualizations to ontological algebras enacts a fundamental epistemological progression. We commence with the mathematical bedrock---sets and relations---devoid of intensional content. We enrich this structure through conceptualization, obtaining abstract domain representations that transcend particular extensional instantiations. We commit our formal languages to these conceptualizations, determining intended models. Finally, we act algebraically upon ontologies, enabling expansion while preserving structural coherence.

At each stage, we confront constitutive limitations:
\begin{itemize}
    \item Extensional structures cannot distinguish essential from accidental relations
    \item Conceptualizations, while intensionally rich, resist direct linguistic expression
    \item Ontological commitments determine intended models but cannot uniquely identify them
    \item Ontologies approximate intended models but cannot fully capture conceptualizations
    \item Expansions generate new structure but remain constrained by algebraic coherence
\end{itemize}

Yet these limitations are not mere deficiencies but rather constitutive features of formal representation itself. The ontology, as explicit specification, becomes the site where language meets conceptual content, where formal rigor confronts the irreducible richness of understanding. It is approximation not by failure but by necessity---the necessary incompleteness of any finite axiomatization before the infinite space of conceptualization.

The ontological algebra emerges from this dialectic as a reconciling structure: it acknowledges the gap between formal specification and conceptual content while providing mechanisms for systematic evolution. The Heyting algebra, with its constructive character, reflects the epistemological situation of ontological work---one cannot assert without construction, cannot negate without evidence, cannot expand without coherence.

%=============================================================================
\section{Summary}
%=============================================================================

The framework developed here comprises the following hierarchy:

\begin{center}
\begin{tikzcd}[row sep=large, column sep=small]
\text{Graded Relational Space } (E_n)_{n \in \mathbb{Z}_+} \\
\text{Multi-Section } s : \mathbb{Z}_+ \to \pfin(E) \arrow[d, "\text{instantiate}"] \\
\text{System } (\domain, S) \arrow[r, leftrightarrow, "\text{vary over } W"] & \text{World } (\domain, w) \\
\text{Phenomenal Cross-Section } \sigma \sqsubseteq S \arrow[d, "\text{flatten}"] \\
\text{Extensional Structure } (\domain, |\sigma|) \arrow[d] \\
\text{Conceptualization } (\domain, \Sigma, \intrel) \arrow[d] \arrow[r] & I(\concept) \\
\text{Ontological Commitment } (\concept, \interp) \arrow[d] \\
\text{Ordered Ontology } (O, \leq) \arrow[d] \\
\text{Ontological Algebra } (\onto, A, \sem{\cdot}) \arrow[d] \\
\text{Expanded Ontology } (O', \leq', \sem{\cdot}')
\end{tikzcd}
\end{center}

The essential contributions:
\begin{enumerate}
    \item The graded relational space and multi-section framework makes explicit the arity-stratification of relational structures, enabling a clean definition of \emph{phenomenal cross-section} as partial observation
    \item The notions of system and world are formally distinguished yet mathematically unified as multi-sections, with their difference residing in interpretive role
    \item A minimal ordered structure---the partial order---suffices for algebraic action, dispensing with topological apparatus
    \item The ontological algebra $(\onto, A, \sem{\cdot})$ bridges static representation and dynamic evolution
    \item Expansion operators provide controlled mechanisms for ontological growth
    \item The Heyting algebra structure ensures constructive, coherent reasoning throughout
\end{enumerate}

Within the meanders between formal precision and conceptual adequacy, between what can be axiomatized and what transcends axiomatization, we discover the living practice of ontological engineering -- an ongoing negotiation between the extensional clarity of logic and the intensional depth of understanding. 

And, even more importantly, through the acknowledgment of a richer structure such as the one above proposed, there emerges a door to an even richer ontological system, one that allows ontologies be raise above being merely a convenient way to rigorously describe how existent entities stand in relation to each other, but one that also \textit{permits} said ontologies to \textit{evolve}, through the rigorous inclusion of ever-new-formed entities, or concepts, by means of a well-defined, precised notion of ``concept/entity derivability".

%=============================================================================
% References
%=============================================================================

\bibliographystyle{plainnat}
\bibliography{formal_ontologies}

\end{document}
